\clearpage\thispagestyle{empty}
\pdfbookmark[0]{Quellen zu Gauß’ Meridianbeobachtungen der Iris und des Neptun}{irisneptune611-sources}
\begin{center}\textbf{Quellen zu Gauß’ Meridianbeobachtungen der Iris und des Neptun}\end{center}
\noindent Diese KI-Übertragung ins heutige Deutsch folgt dem Auszug auf S.~474 der \href{https://gdz.sub.uni-goettingen.de/id/PPN236007467}{\textit{Werke}, Band~VI} (Göttingen, 1874). Carl Friedrich Gauß beobachtete Iris am 21. August 1847 und Neptun am 17. August. Heinrich Christian Schumacher berichtet über die Beobachtungen in den \textit{Mittheilungen} der \textit{Astronomischen Nachrichten}, Band~XXVI, Nr.~611, \href{https://api.digitale-sammlungen.de/iiif/image/v2/bsb10049451_00095/full/1800,/0/default.jpg}{Sp.~173--174}. Die \href{https://api.digitale-sammlungen.de/iiif/image/v2/bsb10049451_00096/full/1800,/0/default.jpg}{Schlusszeile auf Sp.~176} nennt den 16. September 1847. Der Katalog der Bayerischen Staatsbibliothek datiert den gebundenen Band auf 1848.
\par\smallskip\noindent Der Auszug der Werkausgabe hat eine eingeklammerte Überschrift und gibt beide Beobachtungen in seiner eigenen Schreibweise wieder. Er enthält auch die Bemerkung über die Fäden, die vorbeiziehenden Wolken und die nicht beobachtete Kulmination. KI-Übertragung und KI-Satz.
\clearpage
